\documentclass{jsarticle}
\usepackage[dvipdfmx, final]{graphicx}
\usepackage{amsmath}
\usepackage{amssymb}
\usepackage{chemfig}
\begin{document}
\title{}
\author{}
\date{}
\maketitle
 $$  \chemfig{CH_3-CH(-[2]CH_3)-CH_2-C(=[:60]O)-[:-60]O-H} $$ 

 富山大学、1998年度の入学試験の有機化学の問題。
	$$ \chemfig{CH_3-C(=[:90]O)-0H} $$
	ある物質の構造式と、質量を知りたいので、
	原材料の640mgを、分析に掛けて見ると、酢酸が分子量が33mgであり、
	分析で無くなった質量が無いのと、収率の理論値と実存値の
	割合を求める問題である。その上に、誤差が20mgであり、
	分子量を求めると、33mgであり、実材料は、640mgなのに、
	理論値が660mgになっていて、
	二量体である結果、収率は、$ L = {\text{実存量} \over 
	\text{理論値}} $ 
	$$ {330 \over 640} \le 1 $$
	単体量では、誤差が大きすぎることと、
	二量体にすると、
	$$ = {{640} \over {330 \times 2}} \le 1 $$
	であり、収率は、標準では、1以下であることと、誤差が実存量から、
	水素結合になっていることから、
	二量体である。　
	収率を求めると、二量体が適切な解になっている。
	　
  2002年は、ノーベル化学賞に分析化学の田中耕一さんが、誤差のスペクトルを
	見抜くことを、この大学は、同じことをしていた。

\end{document}


